\section{Complete reconstruction-departure evaluation}
\label{app:reconstruction-results}

This appendix specifies the observation laws, fitting procedure, complete
comparison set, stability checks, and resource measurements for
Section~\ref{sec:experiment}. The reproduction materials
contain the implementation, saved predictions and model states, and
machine-readable analyses under
\path{extensions/dimension_forecasting}. All tables and figures in
this appendix are generated from those analyses.

\subsection{Observation laws and data roles}

The linear negative control is
\[
 y_t=0.65y_{t-1}-0.15y_{t-2}+\epsilon_t,
 \qquad \epsilon_t\stackrel{\mathrm{iid}}\sim\mathcal N(0,0.1^2),
\]
where $\mathcal N(\mu,\sigma^2)$ denotes a normal distribution with
mean $\mu$ and variance $\sigma^2$. The process starts with two zeros;
256 stored values are discarded, and 72 values beginning at index 256
are retained. The nonlinear autoregression is
\[
 y_t=\tanh\{1.25y_{t-1}-0.65y_{t-2}
             +0.25\sin(3y_{t-1})+u_t\},
 \qquad u_t\stackrel{\mathrm{iid}}\sim U[-0.1,0.1],
\]
where $U[a,b]$ denotes the uniform distribution on $[a,b]$. Its two
initial values are independent draws from $U[-0.8,0.8]$, and 80 stored
values are discarded before retaining an episode.

The Van der Pol dynamics are $\dot x=v$ and
$\dot v=(1-x^2)v-x$. The Duffing dynamics are $\dot x=v$ and
$\dot v=-0.2v+x-x^3+0.3\cos(1.2t+\varphi)$. Here $x$ is position,
$v$ is velocity, a dot denotes differentiation with respect to continuous
time $t$, and the Duffing phase $\varphi$ is uniform on $[0,2\pi]$.
Initial position and velocity are independent draws from $U[-2,2]$.
The ordinary differential equations are integrated in float64 arithmetic
using the classical fourth-order Runge--Kutta method with step 0.025,
observation interval 0.1, and a discarded initial time interval of length
ten. The scalar observation is $0.95\tanh(x)+\xi$, with independent
sensor noise $\xi\sim U[-0.02,0.02]$. The resulting numerical
integration and observation map define the tested simulation law; the
observations are not exact continuous-time trajectories.

For each law and repetition, the episode sets A, B, C, and D are generated
once. Additional noise has the form $\eta s_A Z$, where
$\eta\in\{0,0.1\}$, $Z\sim\mathcal N(0,1)$ independently at each
stored observation, and $s_A$ is the empirical standard deviation, with
divisor equal to the number of observations, of
all A episode values before this addition. The underlying episodes and
$Z$ arrays are paired across the two noise levels. Added noise affects
both histories and observed future targets. At $\eta=0$, the oscillator
sensor noise is still present. The B/C/D linear observations are Gaussian
conditional on the A-estimated noise amplitude, while their marginal law
across repetitions is a variance mixture. Finite burn-in does not imply
exact stationarity, and the added Gaussian noise means that the observed
responses are not uniformly bounded.

The master seed is 2026110103 and repetitions are indexed by
$r=0,\ldots,23$. A uint32 stream seed is the first output of NumPy's
\texttt{SeedSequence} applied to an integer tuple. Law indices follow
linear AR, nonlinear AR, Van der Pol, and Duffing; set indices are
A/B/C/D $=0/1/2/3$. Base, noise, and retained-row streams use
$(2026110103,r,\mathrm{law},\mathrm{set},\mathrm{purpose})$, with
purposes 1001, 1002, and 1003. Fitting uses
$(2026110103,r,\mathrm{law},1004)$, and the repetition identifier uses
$(2026110103,r,1000)$. The 1,272 intended-distinct stream identifiers
are checked for collisions and for overlap with the 2,173 distinct
identifiers recorded for earlier principal and dimension studies.
Base episodes and standardized noise arrays for different sets use
independent streams. The additional-noise amplitude is estimated on A
and shared across sets. Fitting seeds, row choices, references, and noise
draws are paired across caps and noise levels. Using zero-based indices, the
13 origins $11,15,\ldots,59$ each have a length-12 history ending at the
origin and targets at offsets $1$, $4$, and $12$. A/B/C/D contain
64/32/32/128 episodes. The first three sets retain 512/256/256 histories;
D retains all 1,664 histories. This gives \DFTargetCells{} scored target
values per method over the full experiment. Inference nevertheless uses
24 independent repetition aggregates.

A fits every scaler, reference bank, kernel factor, baseline, and decoder.
B selects penalties and raw-model settings separately by horizon, then
one dimension cap across horizons for each geometric procedure. Its
cap objective uses the uncalibrated predictions in
Equation~\eqref{eq:dimension-cap-selection}; all five cap choices are
saved before C is loaded. C selects each candidate's mixing coefficients
and then one complete non-KAHM procedure within each comparison context.
All fitted states and the complete task roster are sealed before D
episodes are generated. D histories are supplied to fixed predictors;
their targets remain separate until every prediction is sealed. A later
origin can legitimately contain observations that were targets for an
earlier origin in the same episode, but it causes no fitting or
cross-query adaptation. All uncertainty calculations use independent
repetitions, not the overlapping origins.

\subsection{Reference geometry and numerical evaluation}

The geometric bank first standardizes each history coordinate using its
A mean and standard deviation. A standard deviation at most $10^{-12}$
is replaced by one. At most 512 A rows enter the bank fit. Separate
$k$-means partitions request 2, 4, and 8 groups, using five initializations
and at most 150 iterations. If there are too few distinct input rows,
the requested group counts are reduced to the available number. An
unresolved empty-group configuration is retried with halved group counts.
Within each group, at most 24 reference rows are sampled without
replacement. The references are stored in float32. The native PCA converts
its input to float64 and uses the centered sample covariance. Native
candidate cap $c\in\{1,2,4,8,12,20\}$ initializes the requested rank at
$\min\{c,d,n_j-1\}$. Eigenvalues must exceed double-precision machine
epsilon. The requested rank is decremented and the native PCA recomputed
if the smallest range of a projected reference coordinate is below
$10^{-3}$, stopping when the range condition holds or requested rank one
is reached. A rank-one model may therefore remain below that threshold.
The explicit fixed-cap-two comparison instead initializes
$\min\{2,d-1,n_j-1\}$ and uses the same subsequent numerical rules.
The actual stored rank is the number of columns of $P_j$, not the supplied
cap. The forecasting wrapper preserves its explicit zero-rank affine
fallback for singleton or constant references; this is not claimed as
successful native PCA. All cap requests, initial ranks, stored ranks,
reference counts, and reduction or fallback reasons are recorded.

For the row-stored references $Q_j$, define
$\mu_j=n_j^{-1}\sum_{i=1}^{n_j}q_{ji}$ and
\[
 s_j=\max\left\{\frac1{n_j}\sum_{i=1}^{n_j}
                \|q_{ji}-\mu_j\|_2^2,10^{-6}\right\}.
\]
This is the scale used in the common locality weights. The projected
sample covariance is
\[
 S_j=\frac1{n_j-1}\sum_{i=1}^{n_j}
       P_j^\top(q_{ji}-\mu_j)(q_{ji}-\mu_j)^\top P_j.
\]
Its inverse defines the shared Gaussian kernel in
Section~\ref{sec:reconstruction-mechanism}. Kernels and references are
identical between KAHM and positive-RBF reconstruction within a bank.
Independent cap choices can select different projected kernels and bases.

The KAHM regularizer is fitted from reference reconstruction, without
forecast targets. Suppress the local index $j$, let $Q\in\R^{n\times d}$
contain the local references, and define the reference energy
$\nu=\|Q\|_F^2/(nd)$, where $\|\cdot\|_F$ is the square root of the
sum of squared matrix entries. Initialize $\lambda^{(0)}=2.5\nu$.
At iteration $a$, compute
\[
 C^{(a)}=(K+\lambda^{(a)}I_n)^{-1},\qquad
 \lambda^{(a+1)}=2\nu+
     \frac{\|Q-KC^{(a)}Q\|_F^2}{nd}.
\]
The iteration stops when
$|\lambda^{(a+1)}-\lambda^{(a)}|\le10^{-3}$ or after 100 inverse
computations. The reconstruction uses the last computed $C^{(a)}$ and
its corresponding $\lambda^{(a)}$, rather than the subsequent proposal
$\lambda^{(a+1)}$. An unexpected factorization failure or invalid
regularizer stops the experiment; it is not converted into a successful
candidate. This rule determines a reference-based
regularizer; it does not optimize forecasting loss.

For each query, the maximum log-kernel entry is subtracted before
exponentiation. The resulting common positive rescaling leaves both
normalized reconstruction ratios unchanged. Signed cancellation is
checked using $\rho_j$ in Equation~\eqref{eq:mechanism-relative-denominator}.
The affine fallback makes the corresponding raw departure zero if the ratio
cannot be evaluated reliably. The experimental results below report how
often this occurs.

\subsection{Decoders, comparison grids, and descriptive losses}

The common and departure feature columns are standardized separately
using A, with each standard deviation replaced by its maximum with
$10^{-6}$. The blocks are then radially bounded as specified in
Equation~\eqref{eq:mechanism-main-radial}, with no subsequent scaling.
The two unbounded-departure variants retain the bounded common block
and omit only the departure's radial transformation.

Every ridge head uses mean squared error plus
$\alpha\|w\|_2^2$ and an unpenalized intercept, with
$\alpha\in\{10^{-4},10^{-3},10^{-2},10^{-1},1\}$. Equivalently,
the unnormalized sum-of-squares objective uses penalty $n\alpha$ for
$n$ fitting rows. Raw-history ridge and RBF methods standardize their
input coordinates on A with a $10^{-6}$ scale floor. The raw RBF method
concatenates standardized histories with at most 64 Nystr\"om features,
then standardizes the concatenated columns using A. The squared Gaussian
bandwidth equals the median of all strictly positive squared landmark
distances multiplied by $0.5$, $1$, or $2$, with no bandwidth floor.
When there are no strictly positive landmark distances, the median
used by this rule is set to one. B
selects this factor and the ridge penalty separately for each horizon.
For a landmark Gram matrix whose largest eigenvalue is $\zeta_{\max}$,
the Nystr\"om transform retains eigenvalues above
$\max\{10^{-12},10^{-10}\zeta_{\max}\}$.
Let $K_{MM}$ be the symmetrized landmark Gram matrix. Let the columns
of $V$ be its retained orthonormal eigenvectors, and let the diagonal
matrix $\Lambda$ contain their positive eigenvalues. If $k_M(u)$ is the vector of Gaussian similarities between
a standardized query $u$ and the $M$ landmarks, its Nystr\"om coordinates
are $\Lambda^{-1/2}V^\top k_M(u)$. The Gaussian is
$\exp\{-\|u-u'\|_2^2/(2\ell^2)\}$, with squared bandwidth $\ell^2$
chosen by the median-distance rule above. These coordinates are
concatenated with the standardized raw history before column scaling
and ridge fitting.

Histogram gradient boosting uses 100 iterations, L2 regularization
one, and no early stopping. B selects learning rate $0.05$ or $0.1$ and maximum leaf count
7 or 15 separately by horizon. Persistence predicts the latest observed
history value at every horizon. Every non-ridge candidate is then blended
with ridge using a C-selected coefficient from $\{0,1/4,1/2,1\}$.
Ties are resolved by the recorded candidate order. The non-KAHM policy
chooses in this order among ridge, raw RBF, boosting, persistence, the common basis,
and the two RBF-departure variants; both KAHM variants are excluded.
Table~\ref{tab:mechanism-all-methods} reports all methods and conditions.

\begin{table}[!htbp]
\centering\scriptsize
\setlength{\tabcolsep}{3pt}
\caption{All ten procedures, every declared law and both added-noise levels.
Entries are mean horizon-averaged A-variance-normalized D MSE over 24
repetitions; lower is better. The first six columns enter the nonlinear
aggregate with equal weight. The final two form the separate linear control.
These conditionwise summaries are descriptive. ``Unbounded'' omits only the
departure radial bound. Persistence and raw nonlinear controls include their
C-selected blend with ridge.}
\label{tab:mechanism-all-methods}
\begin{tabular}{@{}lrrrrrrrr@{}}
\toprule
Procedure & \multicolumn{2}{c}{Nonlinear AR} & \multicolumn{2}{c}{Van der Pol} & \multicolumn{2}{c}{Duffing} & \multicolumn{2}{c}{Linear AR} \\
 & $\eta=0$ & $0.1$ & $\eta=0$ & $0.1$ & $\eta=0$ & $0.1$ & $\eta=0$ & $0.1$ \\
\midrule
Raw ridge & 0.02200 & 0.03451 & 0.07470 & 0.13584 & 0.01955 & 0.06873 & 0.90079 & 0.90253 \\
Raw RBF & 0.01996 & 0.03243 & 0.01326 & 0.06486 & 0.01296 & 0.04511 & 0.90383 & 0.90556 \\
Boosted trees & 0.02038 & 0.03231 & 0.00380 & 0.05863 & 0.01351 & 0.04908 & 0.90422 & 0.90473 \\
Persistence & 0.02205 & 0.03451 & 0.07470 & 0.13588 & 0.01975 & 0.06955 & 0.90140 & 0.90382 \\
Common features & 0.01928 & 0.03115 & 0.02126 & 0.05902 & 0.01406 & 0.04587 & 0.90139 & 0.90319 \\
Bounded KAHM departure & 0.01928 & 0.03109 & 0.00935 & 0.05735 & 0.01147 & 0.04494 & 0.90230 & 0.90454 \\
Bounded RBF departure & 0.01930 & 0.03119 & 0.01064 & 0.05640 & 0.01289 & 0.04431 & 0.90198 & 0.90366 \\
Unbounded KAHM departure & 0.01949 & 0.03149 & 0.01013 & 0.05929 & 0.01097 & 0.04511 & 0.90237 & 0.90459 \\
Unbounded RBF departure & 0.01937 & 0.03132 & 0.01104 & 0.05861 & 0.01157 & 0.04471 & 0.90244 & 0.90385 \\
Selected non-KAHM & 0.01923 & 0.03110 & 0.00380 & 0.05615 & 0.01175 & 0.04463 & 0.90600 & 0.90700 \\
\bottomrule
\end{tabular}

\end{table}

The machine-readable files \texttt{horizon\_losses.csv} and
\texttt{condition\_horizon\_summary.csv} retain every method, law, noise
level and horizon, including raw MSE and the A-only loss scale. They are not
pooled with historical losses that use different normalizations. The
\texttt{primary\_repetition\_contrasts.csv} file contains all 24 margin-gap values for
each primary contrast, allowing the six lower bounds to be recomputed without
importing any forecasting model.

\begin{table}[!htbp]
\centering\small
\caption{Nonlinear horizon-specific SMSE, averaged equally over the six
law--noise conditions and 24 repetitions. Except for the first two fixed
KAHM rows, procedures use their independently selected caps where applicable.
Every non-ridge forecast includes its C calibration. These descriptive
summaries retain the original A-only variance scales.}
\label{tab:dimension-horizons}
\begin{tabular}{@{}lrrr@{}}
\toprule
Procedure & Horizon 1 & Horizon 4 & Horizon 12 \\
\midrule
KAHM, fixed cap 2 & 0.011654 & 0.021471 & 0.056529 \\
KAHM, native cap 20 & 0.011207 & 0.019108 & 0.066862 \\
Raw ridge & 0.013041 & 0.036954 & 0.127674 \\
Raw RBF & 0.011248 & 0.019724 & 0.063314 \\
Boosted trees & 0.011742 & 0.020626 & 0.056490 \\
Persistence & 0.013056 & 0.036971 & 0.128193 \\
Common features & 0.011149 & 0.018970 & 0.065207 \\
Bounded KAHM departure & 0.011201 & 0.019491 & 0.056043 \\
Bounded RBF departure & 0.011204 & 0.019516 & 0.056647 \\
Unbounded KAHM departure & 0.011378 & 0.020254 & 0.056611 \\
Unbounded RBF departure & 0.011226 & 0.019790 & 0.057292 \\
Selected non-KAHM & 0.011281 & 0.019431 & 0.052623 \\
\bottomrule
\end{tabular}

\end{table}

\subsection{Dimension selection and realized geometry}

Table~\ref{tab:dimension-caps} reports all five geometric procedures'
dimension choices. Native caps twelve and twenty have identical arrays
and forecasts at ambient dimension twelve, and exact tie resolution selects
twelve. Native cap two and the fixed-cap-two policy are also identical.
These identities are checked on bases, features, fitted heads, B choices,
C calibration and final predictions; the policy metadata remain distinct.
Raw ridge, raw RBF, boosting, and persistence are fitted once per task and
shared across the three contexts. The cap-dependent heads and seven banks
are fitted separately without outcome-dependent deduplication.

\begin{table}[!htbp]
\centering\small
\caption{B-selected cap counts. Each procedure contributes 144 nonlinear
and 48 linear tasks; it selects one cap across all three horizons before
C calibration. Exact cap-twelve/cap-twenty ties explain the zero final
column. The common and RBF rows use their own choices, not KAHM's choices.}
\label{tab:dimension-caps}
\begin{tabular}{@{}llrrrrrr@{}}
\toprule
Procedure & Panel & 1 & 2 & 4 & 8 & 12 & 20 \\
\midrule
Common features & Nonlinear & 13 & 10 & 68 & 40 & 13 & 0 \\
Common features & Linear & 35 & 9 & 3 & 0 & 1 & 0 \\
Bounded KAHM departure & Nonlinear & 15 & 22 & 68 & 31 & 8 & 0 \\
Bounded KAHM departure & Linear & 4 & 14 & 18 & 9 & 3 & 0 \\
Bounded RBF departure & Nonlinear & 26 & 29 & 58 & 26 & 5 & 0 \\
Bounded RBF departure & Linear & 14 & 24 & 3 & 5 & 2 & 0 \\
Unbounded KAHM departure & Nonlinear & 17 & 27 & 67 & 25 & 8 & 0 \\
Unbounded KAHM departure & Linear & 2 & 5 & 27 & 11 & 3 & 0 \\
Unbounded RBF departure & Nonlinear & 19 & 39 & 59 & 20 & 7 & 0 \\
Unbounded RBF departure & Linear & 9 & 26 & 10 & 2 & 1 & 0 \\
\bottomrule
\end{tabular}

\end{table}

There are 192 banks and 2,688 local models per cap rule. Local reference
counts range from 2 to 24; 2,670 local reference sets per rule contain 24
rows. Table~\ref{tab:dimension-realized} reports stored ranks and feature
widths. For native cap twenty, 2,685 local models have rank twelve and the
remaining three have two references and rank one. All banks have 14 local
models and a stacked principal-basis rank of twelve. No additional
positive-eigenvalue or projected-range reduction occurs in this sample.
The observed ranks describe these banks, not a guarantee for arbitrary
datasets.

\begin{table}[!htbp]
\centering\small
\setlength{\tabcolsep}{4pt}
\caption{Realized geometry over the complete task roster. Rank statistics
refer to the 2,688 local models in each row; ``Rank 12'' counts full-ambient
local bases. Width ranges refer to 192 complete banks. The full departure
readout adds the common and departure widths. Stacked-basis rank is twelve
for every bank, even when its individual local ranks are small.}
\label{tab:dimension-realized}
\begin{tabular}{@{}lrrrrr@{}}
\toprule
Bank rule & Rank range & Mean rank & Rank 12 & Common width & Departure width \\
\midrule
Fixed cap 2 & 1--2 & 1.999 & 0 & 40--42 & 168 \\
Native cap 1 & 1--1 & 1.000 & 0 & 28--28 & 168 \\
Native cap 2 & 1--2 & 1.999 & 0 & 40--42 & 168 \\
Native cap 4 & 1--4 & 3.997 & 0 & 64--70 & 168 \\
Native cap 8 & 1--8 & 7.992 & 0 & 112--126 & 168 \\
Native cap 12 & 1--12 & 11.988 & 2685 & 160--182 & 168 \\
Native cap 20 & 1--12 & 11.988 & 2685 & 160--182 & 168 \\
\bottomrule
\end{tabular}

\end{table}

The common and departure contributions in Section~\ref{sec:experiment}
use the KAHM-selected bank, not the independently selected common procedure.
The matched common mean loss is 0.032702, compared with 0.031775 for
independently selected common features. Table~\ref{tab:dimension-decomposition}
reports the exact mean loss decomposition relative to fixed cap two.
For native cap twenty, the incremental departure benefit is only
$-0.000653$, compared with $-0.004967$ at cap two. A change in complete
KAHM accuracy must therefore be distinguished from a change in the
incremental value of its reconstruction departure.

\begin{table}[!htbp]
\centering\small
\setlength{\tabcolsep}{4pt}
\caption{Nonlinear loss decomposition relative to the fixed-cap-two
procedure. Total change equals common change plus increment change.
The final column gives the new KAHM-minus-common loss on the same bank.
Negative values in the first two columns denote improvement; a positive
increment change denotes a weaker departure benefit. These are descriptive
identities between calibrated predictor losses.}
\label{tab:dimension-decomposition}
\begin{tabular}{@{}lrrrr@{}}
\toprule
Bank rule & Total change & Common change & Increment change & $L_K-L_0$ \\
\midrule
KAHM-selected bank & -0.000973 & -0.002150 & 0.001177 & -0.003790 \\
Native cap 20 & 0.002507 & -0.001806 & 0.004313 & -0.000653 \\
\bottomrule
\end{tabular}

\end{table}

\subsection{Stability, numerical safeguards and correction activation}

The bounded-feature checks cover 192 tasks, seven banks, all four data
stages, and each of the three bounded procedures. Coefficient checks cover
all five geometric heads and three horizons for each task and bank;
global correction checks apply to the three bounded heads. The
independent analyzer checks the recorded fitted-state bound formulas and
recomputes actual D corrections as the saved forecast minus the saved raw
ridge forecast. It does not claim to reconstruct the fitted coefficient
matrices from summary statistics alone. Table~\ref{tab:mechanism-bounds}
reports maxima over all seven banks and the complete roster; maxima in different columns need
not come from the same task or horizon.

\begin{table}[!htbp]
\centering\small
\setlength{\tabcolsep}{5pt}
\caption{Observed ratios to the deterministic bounded-feature and
ridge-readout limits. Feature ratios concern each procedure's complete
feature vector and its actual radius, including both blocks for the two
appended procedures.
A zero correction with a zero bound contributes ratio zero. ``Violations''
sums failed feature, coefficient and correction checks at the fixed numerical
tolerances. These are implementation checks, not confidence intervals for
forecast risk.}
\label{tab:mechanism-bounds}
\begin{tabular}{@{}lrrrr@{}}
\toprule
Procedure & Feature ratio & Coefficient ratio & Correction ratio & Violations \\
\midrule
Common features & 0.9935 & 0.7154 & 0.2824 & 0 \\
Bounded KAHM departure & 0.9983 & 0.6735 & 0.2207 & 0 \\
Bounded RBF departure & 0.9974 & 0.6738 & 0.2347 & 0 \\
\bottomrule
\end{tabular}

\end{table}

The maximum taskwise ratio of KAHM maximum absolute error to ridge maximum
absolute error is \DFMaxErrorRatio{}, and no task fails a declared stability
screen. All 5,376 task--bank--stage diagnostic records report zero affected
rows for coordinate, gate, kernel, KAHM, and RBF query safeguards. No local
model uses a rank-zero fallback, and no unexpected factor failure occurs.
Counts refer to unions over local models within a task, bank and stage;
the same query would count again if a different bank or stage invoked a
guard. No method fails completely and no final prediction cell requires
recovery. These checks cover every specified task and cap candidate.

\begin{table}[!htbp]
\centering\small
\setlength{\tabcolsep}{4pt}
\caption{C-selected correction activation. Nonlinear and linear counts have
432 and 144 task/horizon cells, respectively. The four mixing-weight columns
combine all 576 cells. A policy choice of raw ridge is counted as zero
correction; raw ridge itself has no mixing selection. Nonzero mixing is not
an empirical guarantee of D improvement, and no minimum activation frequency
was part of the decision.}
\label{tab:mechanism-activation}
\begin{tabular}{@{}lrrrrrr@{}}
\toprule
Procedure & NL active & Linear active & $\gamma=0$ & $1/4$ & $1/2$ & $1$ \\
\midrule
Raw RBF & 432 & 56 & 88 & 53 & 147 & 288 \\
Boosted trees & 416 & 43 & 117 & 104 & 189 & 166 \\
Persistence & 52 & 20 & 504 & 71 & 1 & 0 \\
Common features & 431 & 51 & 94 & 14 & 52 & 416 \\
Bounded KAHM departure & 431 & 52 & 93 & 21 & 73 & 389 \\
Bounded RBF departure & 430 & 53 & 93 & 21 & 62 & 400 \\
Unbounded KAHM departure & 432 & 51 & 93 & 26 & 98 & 359 \\
Unbounded RBF departure & 432 & 48 & 96 & 23 & 68 & 389 \\
Selected non-KAHM & 431 & 72 & 73 & 38 & 99 & 366 \\
\bottomrule
\end{tabular}

\end{table}

\begin{table}[!htbp]
\centering\small
\setlength{\tabcolsep}{8pt}
\caption{Whole-procedure choices made by the C-selected non-KAHM policy.
Each task contributes one choice across all horizons after each candidate's
own horizon-wise C mixing selection. Both KAHM arms are excluded. These are
validation choices, not choices made from D errors.}
\label{tab:mechanism-policy}
\begin{tabular}{@{}lrrrrrrrr@{}}
\toprule
Law & $\eta$ & Ridge & Raw RBF & Boosting & Persist. & Common & RBF & RBF unb. \\
\midrule
Nonlinear AR & 0.0 & 0 & 0 & 0 & 0 & 8 & 3 & 13 \\
Nonlinear AR & 0.1 & 0 & 0 & 0 & 0 & 6 & 6 & 12 \\
Van der Pol & 0.0 & 0 & 0 & 24 & 0 & 0 & 0 & 0 \\
Van der Pol & 0.1 & 0 & 0 & 7 & 0 & 0 & 16 & 1 \\
Duffing & 0.0 & 0 & 2 & 1 & 0 & 0 & 2 & 19 \\
Duffing & 0.1 & 0 & 8 & 2 & 0 & 3 & 6 & 5 \\
Linear AR & 0.0 & 0 & 7 & 5 & 5 & 1 & 3 & 3 \\
Linear AR & 0.1 & 0 & 6 & 6 & 4 & 1 & 1 & 6 \\
\bottomrule
\end{tabular}

\end{table}

\subsection{Measured costs and reproducibility limits}

\begin{table}[!htbp]
\centering\small
\setlength{\tabcolsep}{4pt}
\caption{Mean charged resources over all 192 equal-sized task suites.
Fitting and selection include required A fitting, B selection and C calibration.
Dimension-selected procedures pay for all six native cap candidates.
Shared computation nodes are charged once to every procedure requiring
them; policy fitting pays the union of its non-KAHM candidates. State
denotes the pruned selected deployment object, while arrays count its
owned NumPy buffers. The final two rows give fixed-cap KAHM controls.
These are descriptive implementation costs, not isolated speed benchmarks
or minimum deployment complexity.}
\label{tab:mechanism-resources}
\begin{tabular}{@{}lrrrr@{}}
\toprule
Procedure & Wall (s) & CPU (s) & State (KiB) & Arrays (KiB) \\
\midrule
Raw ridge & 0.001 & 0.001 & 1.1 & 0.6 \\
Raw RBF & 0.015 & 0.015 & 68.1 & 66.7 \\
Boosted trees & 0.681 & 0.681 & 370.5 & 332.1 \\
Persistence & 0.001 & 0.001 & 1.2 & 0.7 \\
Common features & 0.187 & 0.187 & 17.3 & 11.7 \\
Bounded KAHM departure & 0.490 & 0.490 & 136.9 & 129.5 \\
Bounded RBF departure & 0.380 & 0.380 & 68.7 & 61.9 \\
Unbounded KAHM departure & 0.484 & 0.484 & 136.7 & 129.2 \\
Unbounded RBF departure & 0.377 & 0.377 & 68.2 & 61.3 \\
Selected non-KAHM & 1.261 & 1.261 & 136.9 & 123.7 \\
KAHM, fixed cap 2 & 0.085 & 0.086 & 122.1 & 114.7 \\
KAHM, native cap 20 & 0.101 & 0.101 & 183.3 & 175.8 \\
\bottomrule
\end{tabular}

\end{table}

Each task runs in a fresh process with one assigned CPU and one numerical
thread; up to three tasks run concurrently on distinct CPUs. The charged
primitive-node costs exclude diagnostic-only
rank calculations; full-suite timings, including diagnostics, are retained
separately. Node accounting includes model candidates that are actually
evaluated, rather than only the winning hyperparameters. Wall and CPU times
depend on the recorded machine and software environment. The peak whole-task
process RSS is \DFPeakProcessMiB{} MiB; it includes imports, all candidates,
shared diagnostics and evidence arrays, and is not a per-method RAM estimate.
Serialized state and owned array buffers omit transient workspace and do not
measure online update latency. The complete fit stage takes 308.2 seconds
and the prediction stage 227.5 seconds on this shared host, including
worker and evidence overhead. They are not per-method latency estimates.
The study imposes no application-specific cost ceiling.

The analysis program verifies the frozen source and staged input checksums,
the complete task roster, exact target correspondence, A variance scales,
policy predictions, nonfinite/failure accounting and all primary aggregation.
An independent post-release audit regenerates episodes from the pinned
generators and independently checks seeds, noise, windows, targets,
B and C choices, local bases, and all 27 numerical tables. The manuscript
builder verifies the analysis table hashes before producing these displays.

The frozen source inventory has 151 files. Model states are hashed before
and after prediction, all target-bearing files have explicit release
checks, and every candidate forecast is retained. Six disjoint numerical
archives group repetitions 0--3, 4--7, 8--11, 12--15, 16--19, and 20--23;
the first also contains global inventories and seals. Their manifest
records every member checksum. The archived forecast arrays define the
reported scores. Re-evaluating serialized states additionally requires
compatible numerical kernels and package versions; processor-dependent
vectorized arithmetic can change individual floating-point values.

Numerical
development checks, model replay and independent calculations are consistency
checks; they do not constitute a new scientific replication. The protocol was specified locally before evaluation and was informed
by the earlier experiments. It was not publicly preregistered. Neither these known-law repetitions nor the algebraic
distinctness certificate establishes a natural-application forecasting claim.
